% Generated draft: requires mathematical and source review.
\documentclass[11pt]{article}
\usepackage[margin=25mm]{geometry}
\usepackage{fontspec}
\setmainfont{Latin Modern Roman}
\usepackage{amsmath,amssymb,mathtools,unicode-math}
\setmathfont{Latin Modern Math}
\usepackage{enumitem,needspace,xurl,hyperref}
\setlength{\parindent}{0pt}
\setlength{\parskip}{5pt}
\setlength{\emergencystretch}{3em}
\newcommand{\N}{\mathbb{N}}
\newcommand{\Z}{\mathbb{Z}}
\newcommand{\Q}{\mathbb{Q}}
\newcommand{\R}{\mathbb{R}}
\newcommand{\C}{\mathbb{C}}
\newcommand{\F}{\mathbb{F}}
\newcommand{\A}{\mathbb{A}}
\newcommand{\PP}{\mathbb{P}}
\newcommand{\E}{\mathbb{E}}
\newcommand{\eps}{\varepsilon}
\newcommand{\dd}{\,\mathrm{d}}
\begin{document}
\section*{UnsolvedMath ID: 30000068}
\textbf{Source problem code:} catalogue entry
\subsection{Definitions and mathematical statement}
The supplied literature assessment identifies the ambient dimension as three. To spell out the usual fixed-volume relative problem, let \(\Omega\subset\R^3\) be a bounded open convex set and let \(0<v<|\Omega|\). For a measurable set \(E\subset\Omega\), its relative perimeter is
\[
P(E;\Omega)=|D\mathbf1_E|(\Omega)
=\mathcal H^2(\partial^*E\cap\Omega)
\]
when \(E\) has finite perimeter; \(\partial^*E\) denotes its reduced boundary. The standard volume-constrained problem is
\[
I_\Omega(v)=\inf\{P(E;\Omega):E\subset\Omega,\ |E|=v\}.
\]
Only the interface inside \(\Omega\) is charged: the portion of the boundary on \(\partial\Omega\) is free. A minimizer partitions \(\Omega\) into \(E\) and its complement. The phrase ``single disk'' concerns the separating surface, not the three-dimensional volume region \(E\).

On the natural topological reading, the question is whether the closure in \(\overline\Omega\) of the minimizing interior interface is one surface homeomorphic to the closed disk \(\overline{B^2}\), with its boundary on \(\partial\Omega\). This asks for connectedness and disk topology, not necessarily a flat or round Euclidean disk. Whether this is the intended meaning, and the regularity needed to identify the geometric interface with the closure of \(\partial^*E\cap\Omega\), must be checked in the report.

The extracted question omits the volume constraint and its range, boundedness and boundary regularity of \(\Omega\), and whether the intended assertion concerns every minimizer or the existence of a disk-shaped minimizer. Thus the displayed fixed-volume problem supplies a standard interpretation of ``relative isoperimetric partition,'' not a silently completed statement of all missing hypotheses. In particular, without a nontrivial volume or other partition constraint, empty and full sets would trivialize perimeter minimization. The exact source formulation must settle these issues before universal quantifiers over domains, volumes, and minimizers are asserted.

The supplied assessment checked 2026-08-21 describes the single-disk assertion as a conjecture in convex three-dimensional domains and says that the located literature records the conjecture rather than a proof.
\subsection{Short English statement}
The intended question is whether a perimeter-minimizing interface that partitions a convex three-dimensional domain is a single topological disk. The volume constraint, domain regularity, meaning of disk, and quantification over minimizers are missing from the extracted wording.
\subsection{Sources}
\begin{itemize}
\item Local UnsolvedMath registry, permanent ID 30000068, catalogue entry, \emph{Relative Isoperimetric Partitions of Convex Regions}.
\item \emph{Calculus of Variations} (2005), Oberwolfach Reports, pp. 1541--1586, DOI \url{https://doi.org/10.4171/owr/2004/29}; also supplied as ``original Oberwolfach Report 29/2004.''
\item Supplied reference ``recent quantitative relative-isoperimetric results'': \url{https://arxiv.org/abs/2508.21198}.
\item Supplied literature triage checked 2026-08-21; it specifies convex three-dimensional domains but reproduces no exact theorem or full boundary hypotheses.
\end{itemize}
\end{document}
